\documentclass[a4paper, twoside]{report}
%% Language and font encodings
\usepackage[english]{babel}
\usepackage[utf8x]{inputenc}
\usepackage[T1]{fontenc}
\usepackage{textcomp}
\usepackage{caption}
\usepackage{lipsum}
\usepackage{mathtools}
\usepackage{amsmath,amssymb,amsfonts}
\usepackage{rotating}

% ------------------------------------------------------------------------------------------
%% Making a list of acronyms
\usepackage{acro}
\usepackage[acronym]{glossaries}
\makeglossaries

% Make your acronym list here:

\newacronym{ad}{AD}{Alzheimer's Disease} % The argument in the first pair of braces is the reference label, the second is the abbreviation in the text, and the third is the full meaning of the abbreviation.
\newacronym{who}{WHO}{World Health Organization} 

% ------------------------------------------------------------------------------------------



\usepackage[compact]{titlesec}
\titlespacing{\section}{0pt}{2ex}{0ex}
\titlespacing{\subsection}{0pt}{1ex}{1ex}
\titlespacing{\subsubsection}{0pt}{0.5ex}{1ex}

\titleformat{\section}
  {\normalfont\fontsize{18}{40}\bfseries}{\thesection}{1em}{}
  
%% Sets page size and margins
\usepackage[a4paper,top=2.8cm,bottom=2.8cm,left=3cm,right=3cm,marginparwidth=1.75cm]{geometry}

%% Useful packages
\usepackage{amsmath}
\usepackage{graphicx}
\usepackage[colorlinks=true, allcolors=black]{hyperref}

% ------------------------------------------------------------------------------------------
% NOTE: Use this section for the temporary manuscript only
%             Omit this section when printing the final, revised manuscript
%             You should also remove the line: "\linenumbers" below 
%             (before "\input{introduction.tex}")

% Line Numbering on the left side
\usepackage[left]{lineno}
% ------------------------------------------------------------------------------------------


% ------------------------------- TITLE ---------------------------------------

% MODIFY THIS SECTION

\title{Write the title of your thesis here}

% Update names of proponents and supervisor in title.tex

% --------------------------------------------------------------------------------


% ----------------------------------------------------------------------------------
% THIS IS THE BEGINNING OF THE MANUSCRIPT PROPER
% ----------------------------------------------------------------------------------

\begin{document}
\input{title.tex} % Putting the title page

% ------------------------------- ABSTRACT ---------------------------------

% MODIFY THIS SECTION
\section*{Abstract} \large 
\hspace{\parindent} \textit{This is just a sample abstract for demonstration purposes only.} Understanding the brain functions requires the exploration of the activity dynamics and interactions of large neuronal circuits. Two-photon calcium imaging, coupled with the use of genetically encoded calcium indicators (GECIs), is widely used in monitoring large targeted neuronal populations because of its capability to provide single-cell spatial resolution recording in deep brain structures. However, although revolutionary GECIs have sufficient signal-to-noise ratio to resolve single action potentials, there are still several issues that need to be addressed to accurately infer neural activity from raw calcium imaging data. Issues such as neuropil contamination and movement artefacts during imaging of animal models pose problems for the automatic detection and segmentation of neuronal signal sources in two-photon calcium imaging data. In this study, the author presents a standardised computational pipeline that is capable of accurately extracting activity information from automatically segmented cells when performing a comprehensive analysis and exploration of imaging data of any cell model. The whole framework was tested on \textit{in vivo} two-photon calcium imaging data obtained from the mouse hippocampus. Furthermore, the proposed hybrid cell detection and segmentation method was tested on a manually labelled dataset and compared its performance with other existing approaches. It has been shown that \textit{NeuroSEE} has clearly outperformed the other methods in detecting and segmenting cells resulting to a success rate of 75.4\%, as compared to the highest success rate of 67.5\% of other methods. Hence, the proposed pipeline has shown promising results and has demonstrated that efficient artefact correction and the integration of morphology- and activity-based approaches lead to a more accurate extraction of neural information. 

\clearpage
% --------------------------------------------------------------------------------

% -------------------- ACKNOWLEDGEMENTS ------------------------

% MODIFY THIS SECTION

\section*{Acknowledgements} \large

\textbf{\textit{[NOTE]: This is just a sample acknowledgement statement. Kindly modify this accordingly.}}

This work could not have been successfully accomplished without the help of many people.

Thanks to Simon, Bill and Keith for their excellent supervision and support. It has been a pleasure working with them, most especially my primary supervisor, Simon, who constantly gives his support in all of my endeavours throughout my studies. I have definitely enjoyed working with the Schultz Lab, and I look forward to collaborating with them even after graduation.

Thanks to Ann, our postdoc fellow in the lab, for all the help and teaching, especially for the mouse surgery, behavioural training and recording of all the imaging data that I used in my experiments. I have learned a lot from her - both technical and personal. 

Thanks to Yu, our senior research technician in the lab, for always being there to guide me and for being like a second mother to me in the lab. Thanks for keeping me a good company during lunch. 

Thanks to Amanda, our research fellow and co-principal investigator of Simon in the lab, for all the encouragements whenever I give a talk during and outside the lab meeting.

Thanks to all my labmates for warmly welcoming me in the lab and for all the assistance when I was starting. 

Thanks to my closest CDT friends Alessandra, James, Amadeus and De-Shaine for the good friendship and company. I hope we could still keep in touch in the coming years and keep hanging out every time we get 'lost' and feel stressed out with our research.

Thanks to all of my friends in the Philippines and in Hong Kong for keeping in touch with me and for making me feel like I'm still home everyday through our constant messages.

Thanks to my family for their full support on all of my endeavours. Thank you for helping me with my finances, especially during the time that I was having financial difficulties because of scholarship problems. Thanks for always keeping me in touch with my dogs, whom I treat like my own kids, during our video calls.  

Special thanks to my best friends, Josh, JP, Sir Amboy and Grace, for all the unending and unconditional support and for always keeping me sane, especially during my down times. Thanks for believing in me and showing me how proud you are of me. You both are a big part of my success and of who I am today.

And most especially, special thanks to the Lord Almighty for making every bit of my dreams possible. Thanks for all the countless blessings that You have been showering me and for surrounding me with wonderful and loving people. This is all for You!

% --------------------------------------------------------------------------------


% -------------------- OTHER SECTIONS -------------------------------

% DO NOT MODIFY THIS SECTION

% Just modify the contents of the individual sections in their
% respective .tex files

\tableofcontents % Inserts the table of contents
\listoffigures % Inserts the list of figures
\listoftables % Inserts the list of tables


% --------------------------------------------------------------------------------
% Start line numbering
%\linenumbers

% NOTE: Remove this line for the final, revised manuscript

% --------------------------------------------------------------------------------

\input{introduction.tex}
\input{problem.tex}
\input{methods.tex}
%\input{results.tex}
%\input{discussion.tex}
%\input{futurework.tex}
\input{appendix.tex}

\clearpage
\printglossary[type=\acronymtype]


\bibliographystyle{ieeetr}
\bibliography{thesis_refs}

\end{document}